% GENERATED FILE. Edit canonical lessons, then run npm run generate:books.
% canonical-source-hash: fnv1a64:9f66ca148ba500c8
% canonical-lessons: KA-C05-maatanaadu, KA-S05-letter-ta, KA-C05-naanu-kannada-maatanaaduttene, KA-C05-iru, KA-C05-kelasa-maadu, KA-C05-practice

\chapter{The First Verbs}
\label{ch:first-verbs}
\hlchaptermodality{5}

\begin{chapteropening}
\textbf{By the end of this chapter:} \emph{I can build short present-tense Kannada sentences with the chapter's first verbs.}

Say your language, your city and your work in one breath — nānu kannaḍa mātanāḍuttēne, nānu …alli iddēne, nānu kelasa māḍuttēne — and name the three pieces every Kannada verb is stacked from.
\end{chapteropening}

\section[mātanāḍu]{\kn{ಮಾತನಾಡು} (mātanāḍu) — \textquotedblleft{}to speak,\textquotedblright{} your first full verb}

\label{lesson:KA-C05-maatanaadu}



\begin{quote}
So far, \textquotedblleft{}to be.\textquotedblright{} Here is a verb that \emph{does} something — and it shows how Kannada builds \textquotedblleft{}I do X.\textquotedblright{}
\end{quote}

\subsection*{The letters in this word}
\textbf{\kn{ಮಾ}} (\emph{mā}) + \textbf{\kn{ತ}} (\emph{ta}) + \textbf{\kn{ನಾ}} (\emph{nā}) + \textbf{\kn{ಡು}} (\emph{ḍu}) $\to$ \textbf{\kn{ಮಾತನಾಡು}} (\emph{mātanāḍu}).

\begin{cousinweb}
\textbf{\kn{ಮಾತನಾಡು}} (\emph{mātanāḍu}, \textquotedblleft{}to speak, to talk\textquotedblright{}) is native Dravidian, built on \textbf{\kn{ಮಾತು}} (\emph{mātu}, \textquotedblleft{}word, speech\textquotedblright{}) + \emph{āḍu} (\textquotedblleft{}to play/do\textquotedblright{}) — \textquotedblleft{}to word-play,\textquotedblright{} i.e. to converse. To say \textquotedblleft{}\textbf{I speak},\textquotedblright{} Kannada adds a two-part ending to the stem:

\begin{quote}
\emph{mātanāḍ-} + \emph{-utt-} (present) + \emph{-ēne} (\textquotedblleft{}I\textquotedblright{}) $\to$ \textbf{\kn{ಮಾತನಾಡುತ್ತೇನೆ}} (\emph{mātanāḍuttēne}), \textquotedblleft{}I speak.\textquotedblright{}
\end{quote}

The \textbf{-\kn{ಏನೆ}} (\emph{-ēne}) ending \emph{is} \textquotedblleft{}I,\textquotedblright{} so \emph{nānu} is rarely needed. Change the ending, change the person: \emph{mātanāḍuttāne} (he), \emph{mātanāḍuttāḷe} (she), \emph{mātanāḍuttāre} (they / you-respectful).
\end{cousinweb}

\begin{grammarlens}[title={Grammar Lens: stem + tense + person}]
Every Kannada verb is built this way — a stem, a tense-marker, a person-ending. And, like Tamil and Telugu, Kannada marks \textbf{no gender in the first person}: \emph{mātanāḍuttēne} is the same for a man or a woman.
\end{grammarlens}

\subsection*{Your turn}
\begin{itemize}
\raggedright
  \item \emph{Say it:} \textquotedblleft{}mātanāḍu\textquotedblright{} (speak)
  \item \emph{Say it:} \textquotedblleft{}I speak\textquotedblright{} — \emph{mātanāḍuttēne}
  \item \emph{Say it:} the word \textquotedblleft{}speech\textquotedblright{} inside it (\emph{mātu})
  \item \emph{Recall:} say \emph{nānu}, then read \textbf{\kn{ಚೆನ್ನಾಗಿ}}
\end{itemize}

\begin{tcolorbox}[breakable,colback=teal!4,colframe=teal!35!black,title={Before you move on}]
Build \textquotedblleft{}I speak\textquotedblright{} from \emph{mātanāḍu}, and name the noun it's built on. (\emph{Mātanāḍuttēne}; \emph{mātu}, \textquotedblleft{}word / speech.\textquotedblright{})
\end{tcolorbox}
\section[ta]{\kn{ತ} — one character, met inside words you already say}

\label{lesson:KA-S05-letter-ta}



\begin{quote}
Before the new one: \kn{◌ಿ} — what does it do?

One character this time. Just one — and you have been saying it for pages without knowing which mark on the page it was.
\end{quote}

\subsection*{Script you'll notice: \kn{ತ}}
\textbf{\kn{ತ}} — \emph{ta}.

It is a \textbf{consonant}, and in this script a consonant is never bare: it comes with an \emph{a} already in it. So it is not \emph{t}, it is \textbf{ta}.

You already say these, and every one of them has \kn{ತ} somewhere inside it:

\begin{itemize}
\raggedright
  \item \textbf{\kn{ಸಂತೋಷ}} \emph{santōṣa} — joy / pleased to meet you
  \item \textbf{\kn{ಹೋಗಿ} \kn{ಬರುತ್ತೇನೆ}} \emph{hōgi baruttēne} — goodbye (lit. \textquotedblleft{}I'll go and come back\textquotedblright{})
  \item \textbf{\kn{ಮತ್ತೆ} \kn{ಸಿಗೋಣ}} \emph{matte sigōṇa} — we'll meet again
  \item \textbf{\kn{ಮಾತನಾಡು}} \emph{mātanāḍu} — to speak
\end{itemize}

\subsection*{Writing: \kn{ತ} — copy what you see}
Put your pen on \kn{ತ} and follow its line. Copy the shape you can see — slowly, and larger than it is printed.

\begin{quote}
This book does not yet tell you \textbf{where to start the character or which way to travel}. That is a real thing, taught with real variation from school to school, and it is not written down here until it can be written down with a source. Copying what is in front of you needs no such source, and it is how the shape gets into your hand in the meantime.
\end{quote}

\subsection*{Your turn}
\begin{itemize}
\raggedright
  \item \emph{Look:} at these words, and find \kn{ತ} in the ones that have it
\end{itemize}

\begin{quote}
\kn{ಸಂತೋಷ}  ·  \kn{ಹೋಗಿ} \kn{ಬರುತ್ತೇನೆ}  ·  \kn{ನಮಸ್ಕಾರ}
\end{quote}

\begin{itemize}
\raggedright
  \item \emph{Trace it:} \kn{ತ} three times, saying \emph{ta} as you finish each one
  \item \emph{Look:} back at any page of this chapter and find \kn{ತ} once more
\end{itemize}

\begin{tcolorbox}[breakable,colback=teal!4,colframe=teal!35!black,title={Before you move on}]
Which character is this — \kn{ತ}? What sound does it carry? (\emph{\textbf{ta}}.) Name one word you already say that contains it.
\end{tcolorbox}
\section[nānu kannaḍa mātanāḍuttēne]{\kn{ನಾನು} \kn{ಕನ್ನಡ} \kn{ಮಾತನಾಡುತ್ತೇನೆ} (nānu kannaḍa mātanāḍuttēne) — \textquotedblleft{}I speak Kannada\textquotedblright{}}

\label{lesson:KA-C05-naanu-kannada-maatanaaduttene}



\begin{quote}
Your first full, moving sentence — and it names one of the classical languages of India.
\end{quote}

\subsection*{The letters in this word}
\textbf{\kn{ಕನ್ನಡ}} (\emph{kannaḍa}): \textbf{\kn{ಕ}} (\emph{ka}) + \textbf{\kn{ನ್ನ}} (\emph{nna}, doubled) + \textbf{\kn{ಡ}} (\emph{ḍa}).

\begin{cousinweb}
\textbf{\kn{ನಾನು} \kn{ಕನ್ನಡ} \kn{ಮಾತನಾಡುತ್ತೇನೆ}} = \textbf{\kn{ನಾನು}} (\emph{nānu}, \textquotedblleft{}I\textquotedblright{}) + \textbf{\kn{ಕನ್ನಡ}} (\emph{kannaḍa}, the object) + \textbf{\kn{ಮಾತನಾಡುತ್ತೇನೆ}} (\emph{mātanāḍuttēne}, \textquotedblleft{}speak\textquotedblright{}) — \textquotedblleft{}I Kannada speak,\textquotedblright{} the verb last as always. The name \textbf{\kn{ಕನ್ನಡ}} (\emph{kannaḍa}) is native (often linked to \emph{kar-nāḍu}, \textquotedblleft{}the black/high land,\textquotedblright{} whence \emph{Karnataka}). Kannada has a literature over a thousand years old and holds the second-most Jnanpith literary awards of any Indian language — a classical tongue, spoken by some fifty million.
\end{cousinweb}

\begin{grammarlens}[title={Grammar Lens: no gender on \textquotedblleft{}I speak\textquotedblright{}}]
\emph{Mātanāḍuttēne} is the same whether a man or a woman says it — Kannada, like Tamil and Telugu, marks \textbf{no gender in the first or second person} (only the third: \emph{-āne} he, \emph{-āḷe} she). This is unlike Hindi, where even \textquotedblleft{}I speak\textquotedblright{} changes for the speaker's gender (\emph{boltā/boltī}).
\end{grammarlens}

\subsection*{Your turn}
Say these aloud:

\begin{itemize}
\raggedright
  \item \textquotedblleft{}nānu kannaḍa mātanāḍuttēne\textquotedblright{}
  \item the likely meaning of \emph{kannaḍa} (\emph{kar-nāḍu}, \textquotedblleft{}the high/black land\textquotedblright{} $\to$ \emph{Karnataka})
  \item does \textquotedblleft{}I speak\textquotedblright{} change for a man vs. a woman? (No — no 1st-person gender)
\end{itemize}

\begin{tcolorbox}[breakable,colback=teal!4,colframe=teal!35!black,title={Before you move on}]
Say \textquotedblleft{}I speak Kannada,\textquotedblright{} and give the likely root of the name. (\emph{Nānu kannaḍa mātanāḍuttēne}; \emph{kar-nāḍu}, \textquotedblleft{}the high/black land.\textquotedblright{})
\end{tcolorbox}
\section[iru]{\kn{ಇರು} (iru) — \textquotedblleft{}to be, to stay, to live\textquotedblright{}}

\label{lesson:KA-C05-iru}



\begin{quote}
You have already used this verb — in \textquotedblleft{}how are you.\textquotedblright{} Now meet it in its own right, and use it to say where you live.
\end{quote}

\subsection*{The letters in this word}
\textbf{\kn{ಇ}} (independent \emph{i}) + \textbf{\kn{ರು}} (\emph{ru}) $\to$ \textbf{\kn{ಇರು}} (\emph{iru}).

\begin{cousinweb}
\textbf{\kn{ಇರು}} (\emph{iru}, \textquotedblleft{}to be, to exist, to stay, to live\textquotedblright{}) is native Dravidian — the very same \emph{iru} Tamil uses, and the verb behind Chapter 3's \emph{hēgiddīrā} (\textquotedblleft{}how are you\textquotedblright{}) and \emph{cennāgiddēne} (\textquotedblleft{}I'm well\textquotedblright{}). It also means \textquotedblleft{}to reside\textquotedblright{}: \textbf{\kn{ನಾನು} \kn{ಬೆಂಗಳೂರಿನಲ್ಲಿ} \kn{ಇದ್ದೇನೆ}} (\emph{nānu Beṅgaḷūrinalli iddēne}, \textquotedblleft{}I live in Bangalore\textquotedblright{}). One verb covers \textquotedblleft{}be,\textquotedblright{} \textquotedblleft{}stay,\textquotedblright{} and \textquotedblleft{}live.\textquotedblright{}
\end{cousinweb}

\begin{grammarlens}[title={Grammar Lens: \textquotedblleft{}in\textquotedblright{} comes after the noun}]
\textbf{\kn{ಬೆಂಗಳೂರಿನಲ್ಲಿ}} (\emph{Beṅgaḷūrinalli}) = \emph{Bengaḷūru} + \textbf{-\kn{ಅಲ್ಲಿ}} (\emph{-alli}, \textquotedblleft{}in\textquotedblright{}), glued to the \textbf{end} of the noun. Kannada uses \textbf{post}positions (case-suffixes) where English uses prepositions: \textquotedblleft{}Bangalore-in I am,\textquotedblright{} the verb last.
\end{grammarlens}

\subsection*{Your turn}
\begin{itemize}
\raggedright
  \item \emph{Say it:} \textquotedblleft{}iru\textquotedblright{}
  \item \emph{Say it:} \textquotedblleft{}I live in Bangalore\textquotedblright{} — \emph{nānu Beṅgaḷūrinalli iddēne}
  \item \emph{Say it:} where \textquotedblleft{}in\textquotedblright{} (\emph{-alli}) sits — glued to the end of the noun
  \item \emph{Recall:} read \textbf{\kn{ಹೋಗು}}
\end{itemize}

\begin{tcolorbox}[breakable,colback=teal!4,colframe=teal!35!black,title={Before you move on}]
Name three things the one verb \emph{iru} can mean, and say where \textquotedblleft{}in\textquotedblright{} goes. (\textquotedblleft{}To be / to stay / to live\textquotedblright{}; \emph{-alli} \textquotedblleft{}in\textquotedblright{} attaches to the \textbf{end} of the noun.)
\end{tcolorbox}
\section[kelasa māḍu]{\kn{ಕೆಲಸ} \kn{ಮಾಡು} (kelasa māḍu) — \textquotedblleft{}to work\textquotedblright{}}

\label{lesson:KA-C05-kelasa-maadu}



\begin{quote}
The last verb of the chapter — and, like Hindi's \emph{karnā} and Tamil's \emph{sey}, it is a \textquotedblleft{}do\textquotedblright{} verb that builds a hundred others.
\end{quote}

\subsection*{The letters in this word}
\textbf{\kn{ಕೆಲಸ}} (\emph{kelasa}, \textquotedblleft{}work\textquotedblright{}): \textbf{\kn{ಕೆ}} (\emph{ke}) + \textbf{\kn{ಲ}} (\emph{la}) + \textbf{\kn{ಸ}} (\emph{sa}). \textbf{\kn{ಮಾಡು}} (\emph{māḍu}, \textquotedblleft{}do\textquotedblright{}): \textbf{\kn{ಮಾ}} (\emph{mā}) + \textbf{\kn{ಡು}} (\emph{ḍu}).

\begin{cousinweb}
\textbf{\kn{ಮಾಡು}} (\emph{māḍu}, \textquotedblleft{}to do, to make\textquotedblright{}) is a core native Dravidian verb, and Kannada — like Hindi with \emph{karnā}, Tamil with \emph{sey}, and Telugu with \emph{cēyu} — builds compound verbs by putting a noun in front of it:

\begin{itemize}
\raggedright
  \item \textbf{\kn{ಕೆಲಸ} \kn{ಮಾಡು}} (\emph{kelasa māḍu}) = \textquotedblleft{}work-do\textquotedblright{} = \textquotedblleft{}\textbf{to work}\textquotedblright{} (\emph{kelasa}, \textquotedblleft{}work\textquotedblright{})
  \item \textbf{\kn{ಅಡುಗೆ} \kn{ಮಾಡು}} (\emph{aḍuge māḍu}) = \textquotedblleft{}to cook\textquotedblright{}
  \item \textbf{\kn{ಸಹಾಯ} \kn{ಮಾಡು}} (\emph{sahāya māḍu}) = \textquotedblleft{}to help\textquotedblright{} (\emph{sahāya} $\leftarrow$ Sanskrit)
\end{itemize}

So \textquotedblleft{}I work\textquotedblright{} is \textbf{\kn{ನಾನು} \kn{ಕೆಲಸ} \kn{ಮಾಡುತ್ತೇನೆ}} (\emph{nānu kelasa māḍuttēne}). Notice \emph{sahāya māḍu} (\textquotedblleft{}to help\textquotedblright{}) pairs a \textbf{Sanskrit} noun with the \textbf{native} \emph{māḍu} — Kannada's two vocabularies working together in one verb.
\end{cousinweb}

\begin{grammarlens}[title={Grammar Lens: noun + \kn{ಮಾಡು} = a new verb}]
This \textquotedblleft{}noun + \emph{māḍu}\textquotedblright{} pattern is the exact twin of Hindi's \textquotedblleft{}noun + \emph{karnā},\textquotedblright{} Tamil's \textquotedblleft{}noun + \emph{sey},\textquotedblright{} and Telugu's \textquotedblleft{}noun + \emph{cēyu}.\textquotedblright{} Four languages — two of them from an unrelated family — all landed on the same trick.
\end{grammarlens}

\subsection*{Your turn}
Say these aloud:

\begin{itemize}
\raggedright
  \item \textquotedblleft{}kelasa māḍu\textquotedblright{} (to work)
  \item \textquotedblleft{}I work\textquotedblright{} — \emph{nānu kelasa māḍuttēne}
  \item the Hindi, Tamil, and Telugu twins of this trick (\emph{karnā}, \emph{sey}, \emph{cēyu})
\end{itemize}

\begin{tcolorbox}[breakable,colback=teal!4,colframe=teal!35!black,title={Before you move on}]
How does Kannada turn \textquotedblleft{}work\textquotedblright{} into \textquotedblleft{}to work,\textquotedblright{} and what verbs work the same way in its neighbours? (\emph{Kelasa} + \emph{māḍu} = \textquotedblleft{}work-do\textquotedblright{}; Hindi \emph{karnā}, Tamil \emph{sey}, Telugu \emph{cēyu}.)
\end{tcolorbox}
\section[Practice]{Chapter 5 — The first verbs}

\label{lesson:KA-C05-practice}



\begin{quote}
No new words. Three verbs, one engine — build real sentences about yourself.
\end{quote}

\subsection*{Your turn: Build the sentences}
Say these aloud:

\begin{itemize}
\raggedright
  \item \textquotedblleft{}I speak Kannada\textquotedblright{} — \emph{nānu kannaḍa mātanāḍuttēne}
  \item \textquotedblleft{}I live in Bangalore\textquotedblright{} — \emph{nānu Beṅgaḷūrinalli iddēne}
  \item \textquotedblleft{}I work\textquotedblright{} — \emph{nānu kelasa māḍuttēne}
  \item \textquotedblleft{}he speaks\textquotedblright{} vs. \textquotedblleft{}she speaks\textquotedblright{} — \emph{mātanāḍuttāne} / \emph{mātanāḍuttāḷe}
\end{itemize}

\subsection*{You'll want to know: The engine}
Say these aloud:

\begin{itemize}
\raggedright
  \item the three stacked pieces of every Kannada verb (stem + tense + person)
  \item the ending that means \textquotedblleft{}I\textquotedblright{} (\emph{-ēne})
  \item where Kannada marks gender — and where it doesn't (3rd person only)
\end{itemize}

\begin{cousinweb}
\begin{itemize}
\raggedright
  \item \emph{mātanāḍu} \textquotedblleft{}to speak\textquotedblright{} $\leftarrow$ \emph{mātu} \textquotedblleft{}word\textquotedblright{}
  \item \emph{iru} \textquotedblleft{}to be / stay / live\textquotedblright{}; \emph{-alli} \textquotedblleft{}in\textquotedblright{} attaches to the noun's end
  \item \emph{māḍu} \textquotedblleft{}to do\textquotedblright{} $\to$ \emph{kelasa māḍu} \textquotedblleft{}to work\textquotedblright{} (the twin of \emph{karnā} / \emph{sey} / \emph{cēyu})
  \item \emph{kannaḍa} $\leftarrow$ \emph{kar-nāḍu}, \textquotedblleft{}the high/black land\textquotedblright{} ($\to$ \emph{Karnataka})
\end{itemize}
\end{cousinweb}

\begin{tcolorbox}[breakable,colback=teal!4,colframe=teal!35!black,title={Before you move on}]
In one breath, say your language, your city, and your work in Kannada. (\emph{Nānu kannaḍa mātanāḍuttēne. Nānu … alli iddēne. Nānu kelasa māḍuttēne.})
\end{tcolorbox}
